\documentclass[11pt,a4paper]{article}
% Shared preamble for RuPhO English problem statements (match T1 2026 style).
% Use: \input{latex/rupho_en_preamble.tex} after \documentclass.
\usepackage[margin=1in]{geometry}
\usepackage{amsmath,amssymb}
\usepackage{graphicx}
\usepackage{float}
\usepackage{enumitem}
\usepackage{microtype}
\usepackage{caption}
\usepackage{tabularx}
\usepackage{hyperref}

\captionsetup{font=small,labelfont=bf,skip=6pt}

\setlist[itemize]{leftmargin=1.2cm}
\setlist[enumerate]{leftmargin=1.2cm}

% One outer box per subproblem: [id | statement | points] (no inner rules)
\newcommand{\problembox}[3]{%
  \par\addvspace{0.42em}%
  \noindent
  \setlength{\fboxsep}{7.5pt}%
  \setlength{\fboxrule}{0.95pt}%
  \fbox{%
    \begin{minipage}{\dimexpr\linewidth-2\fboxsep-2\fboxrule\relax}
    \begin{tabularx}{\linewidth}{@{}>{\raggedright\arraybackslash\bfseries}p{2.1em} @{\hspace{0.9em}} >{\raggedright\arraybackslash}X @{\hspace{0.9em}} >{\raggedleft\arraybackslash\bfseries}p{2.4em}@{}}
    #1 & #2 & #3
    \end{tabularx}
    \end{minipage}%
  }%
  \par\addvspace{0.32em}%
}

\title{\textbf{Phase Transitions in Solutions}}
\author{\normalsize Y16 (2016) --- English translation}
\date{}
\begin{document}
\maketitle

Liquids can transform into solid and gaseous states via first-order phase transitions (i.e., transitions associated with the absorption or release of latent heat). These transitions occur at a constant temperature. In reference tables, these temperatures are usually provided for standard conditions, though their actual values depend on many factors. In this problem, we consider some of these dependencies (restricting our attention exclusively to first-order phase transitions).

The phase transition curve of a substance from state $1$ to state $2$ is described by the Clapeyron--Clausius equation:
$$\frac{dp}{dT} = \frac{q_{12}}{T(v_2-v_1)}$$
where $p$ is the pressure, $T$ is the temperature, $q_{12}$ is the specific latent heat of the transition $1 \to 2$, and $v_1, v_2$ are the specific volumes of the respective phases.

\problembox{A1}{At what temperature does water boil in Switzerland at Dufour Peak at an altitude of $4634~\text{m}$, if the atmospheric pressure there is $430~\text{mm~Hg}$? Normal atmospheric pressure is $760~\text{mm~Hg}$. Under normal conditions, water boils at $100^{\circ}\text{C}$.}{0.20}

Surface tension $\sigma$ of a liquid can affect the saturated vapor pressure of a substance above a curved surface. Let $K$ be the curvature of the surface, defined as the reciprocal of the radius of curvature ($K = 1/R$).

\problembox{A2}{Find the Laplace pressure in the case of a cylindrical surface (the pressure difference between opposite sides of the curved liquid interface). Express your answer in terms of $\sigma$, $K$, $p_a$ (atmospheric pressure), and physical constants.}{0.10}

\problembox{A3}{Find the Laplace pressure for a spherical surface. Express your answer in terms of $\sigma$, $K$, $p_a$ (atmospheric pressure), and physical constants.}{0.10}

\problembox{A4}{Consider a gas bubble inside a liquid. Where is the pressure greater: inside ($p_{\text{in}}$) or outside ($p_{\text{out}}$) the bubble?}{0.10}

In the next part of the problem, we consider how the curvature of the liquid surface affects its vapor pressure. To analyze this, consider communicating cylindrical vessels placed in a thermostat maintained at a constant temperature $T$ (Fig. 1 illustrates the case where the liquid completely wets the walls of the vessels). The width of one of the vessels is small enough for capillary effects to be observable. The vapor pressure above the flat surface $AB$ is $p_1$, and above the curved surface $CD$ is $p_2$. Let $K$ be the curvature of the surface $CD$. The specific volume of the liquid is $v_L$, and that of the vapor is $v_G$ (assume that the vapor density varies negligibly over the height $h$).

\begin{figure}[H]
    \centering
    \includegraphics[width=0.72\textwidth]{T2_EN_assets/fig1.jpg}
    \caption{Saturated vapor pressure over a curved surface.}
\end{figure}

\problembox{A5}{Express the height difference $h$ in terms of $\sigma$, $K$, $v_L$, $v_G$, $T$, the molar mass $\mu$, and physical constants.}{0.20}

\problembox{A6}{Assuming the liquid wets the walls of the vessel, express the pressure difference $p_2-p_1$ in terms of $\sigma$, $K$, $v_L$, $v_G$, $T$, $\mu$, and physical constants.}{0.40}

\problembox{A7}{What would the answer be if the liquid does not wet the walls of the vessel? Express your answer in terms of $\sigma$, $K$, $v_L$, $v_G$, $T$, $\mu$, and physical constants.}{0.10}

\problembox{A8}{For a highly curved surface ($\sigma K \gg |p_2 - p_1|$), the vapor density can no longer be assumed constant. What will be the pressure difference $p_2-p_1$ in this case? Express your answer in terms of $\sigma$, $K$, $v_L$, $v_G$, $T$, $\mu$, and physical constants.}{0.40}

Next, let us examine how the conditions for phase transitions change when a solute is dissolved in a liquid. First, consider osmosis---the one-way diffusion of solvent molecules through a semi-permeable membrane. This membrane allows only solvent molecules to pass through. Solvent molecules diffuse towards the region of higher solute concentration. Let the number density of the solvent molecules be $n$, and that of the solute molecules be $n_s$. This phenomenon gives rise to an osmotic pressure $p_{\text{os}}$---the excess pressure that must be applied to the solution to halt the net diffusion of solvent molecules. Fig. 2 schematically shows an osmometer, a device used to measure osmotic pressure.

\begin{figure}[H]
    \centering
    \includegraphics[width=0.72\textwidth]{T2_EN_assets/fig2.jpg}
    \caption{An osmometer.}
\end{figure}

\problembox{B1}{Determine the osmotic pressure $p_{\text{os}}$ for a dilute solution (a solution with a low solute concentration). Express your answer in terms of $p_a$, $n$, $n_s$, $T$, and physical constants.}{0.50}

Consider the change in saturated vapor pressure when a small amount of substance is dissolved in a liquid. In Fig. 3, the pure solvent (water) and the solution are separated by a semi-permeable membrane. Let $p_0$ be the saturated vapor pressure above the pure water, and $p$ be the vapor pressure above the solution.

\begin{figure}[H]
    \centering
    \includegraphics[width=0.72\textwidth]{T2_EN_assets/fig3.jpg}
    \caption{Vapor pressure over a solution.}
\end{figure}

\problembox{B2}{Express the pressure difference $p-p_0$ in terms of $v_L$, $v_G$, $p_{\text{os}}$, and physical constants.}{0.50}

\problembox{B3}{Express the pressure difference $p-p_0$ in terms of $p_0$, $n$, $n_s$, and physical constants.}{0.30}

In the final part of the problem, consider three distinct phases of a substance ($1$ --- liquid, $2$ --- vapor, $3$ --- solid). Let $q_{21}$, $q_{23}$, and $q_{13}$ denote the specific latent heats released during the transitions $2\to{1}$, $2\to{3}$, and $1\to{3}$, respectively ($q_{21}>0$, $q_{23}>0$, $q_{13}>0$). Analyze the system near its three-phase coexistence point (the triple point).

\problembox{C1}{Assuming the difference $|p - p_0|$ is small, find the shift in the boiling point of the liquid when a substance is dissolved in it. Express your answer in terms of $\mu$, $T$, $q_{12}$, $n_s$, $n$, and physical constants.}{2.00}

\problembox{C2}{Determine the sign of the change in the boiling point of a liquid when a solute is dissolved in it.}{0.30}

\problembox{C3}{Express $q_{23}$ in terms of $q_{13}$ and $q_{21}$ at the triple point.}{1.50}

\problembox{C4}{Qualitatively sketch the phase equilibrium curves between all three phases near the triple point in $(p, T)$ coordinates. Assume $v_G \gg v_L$ and $v_G \gg v_s$ (where $v_s$ is the specific volume of the solid phase).}{1.00}

\problembox{C5}{Find the shift in the freezing point of the liquid when a substance is dissolved in it. Express your answer in terms of $\mu$, $T$, $q_{13}$, $n_s$, $n$, and physical constants.}{2.00}

\problembox{C6}{Determine the sign of the change in the freezing point of a liquid when a solute is dissolved in it.}{0.30}

\vspace{1em}
\noindent\textit{Source:} \url{https://pho.rs/p/3028}

\end{document}